\scanpage{35}
une \unsure{extension} de la notion de \uline{morphismes de graphes}.

\deleted{\unsure{6}}\qquad\deleted{Soit}

6. \uline{Morphismes de générisation de graphes \deleted{\unsure{quelconques}}\inserted{et de maquettes}}

Soit $G=(S,\widetilde A,\sigma_{\widetilde A},o_{\widetilde A}:\widetilde A\unsure{\to|{-}}S)$ un graphe ordinaire. On lui associe de façon bien connue un \uline{groupoïde fondamental} \inserted{$\Pi_1(G)$}, dont l'ens des objets est $S$, et l'ens des \deleted{\unsure{chemins de}}\inserted{morphismes} de $s$ à $s'$ est l'ens des ``chemins simples'' de $s$ à $s'$, qui correspondent aux classes d'homotopie de chemins de $s$ à $s'$ dans la réalisation topologique du graphe $G$. Le groupoïde $\Pi_1(G)$ dépend fonctoriellement (au sens le plus strict !) du graphe $G$.

D'ailleurs, un morphisme de graphes
\[
(41)\qquad f:G'\longrightarrow G^{\deleted{\unsure{'}}}
\]
\mspage{34}
